\chapter{Введение}
\label{chap:intro}
Протокол языкового сервера (LSP), представленный в 2016 году
\cite{language-server-protocol:2016}, позволяет отделить логику поддержки языка
от редактора кода. Это решает проблему $M \times N$, позволяя одному языковому
серверу работать с множеством редакторов и IDE.

\section{Язык программирования Jasmin}

Криптографическое ПО играет ключевую роль в безопасности, и уязвимости в нём,
могут иметь катастрофические последствия. Для снижения рисков создания ПО с
ошибками создаются фреймворки, позволяющие формально верифицировать
криптографические примитивы.  Одним из таких фреймворков является Jasmin
\cite{jasmin:2017}.

Jasmin — это фреймворк для реализации и верификации криптографических
примитивов. Его компилятор позволяет транслировать код в спецификации EasyCrypt
\cite{easycrypt:2013} для формального анализа. Данная работа посвящена
реализации языкового сервера для Jasmin.

В данной работе мы представляем реализацию инфраструктуры языкового сервера для
языка программирования Jasmin.

Основной вклад этой работы — реализация языкового сервера, который поддерживает:
\begin{itemize}
    \item Диагностику для языка Jasmin
    \item Переход к определению ($\texttt{textDocument/definition}$)
    \item Подсветку ссылок ($\texttt{textDocument/highlight}$)
    \item Отображение типов при наведении ($\texttt{textDocument/hover}$)
    \item Поддержку Tree-sitter для языка Jasmin
\end{itemize}